\section{Introduction}
\label{sec:introduction}
Generative image editing has made it possible to remove objects, replace people, or insert convincing content with only a short instruction. Inpainting methods alter a selected part of a photograph while producing a visually coherent surrounding scene. For image forensics, the task is therefore not only to decide whether an image has been edited, but also to identify which pixels support that decision. The detector considered in this paper receives a single image; the original photograph is not required at inference.

Local edits are difficult because the altered area can occupy only a small fraction of a high-resolution image. Resizing the entire image makes global context available at manageable cost, but may suppress the details of a small manipulation. Inspecting high-resolution crops preserves these details, but a crop may contain no edited pixels, or may lack the context needed to distinguish a legitimate object from generated content. Genuine textures, compression, and camera processing can also resemble manipulation cues. A reliable system must consider both missed edits and false alarms on genuine photographs.

Prior methods study manipulation-sensitive residuals, pretrained visual representations, and their combination for image detection and localization~\cite{trufor,declip,opensdi,fused}. OpenSDID evaluates diffusion-based inpainting across several generators~\cite{opensdi}. FUSED combines semantic and forensic representations and evaluates transfer across generators and held-out benchmarks~\cite{fused}. However, an image-level score alone does not show whether the predicted evidence lies inside the edited region; detectors may exploit traces elsewhere in the image~\cite{inpx}. Measuring localization by edit size and examining errors on genuine images are therefore necessary complements to overall accuracy.

We investigate a detector that combines contextual visual features with residual features and produces both an image-level score and a localization map. The approach processes a resized whole image and overlapping native-resolution crops, bringing global context and local detail into the prediction. We additionally examine joint-view supervision: the whole view receives the image label, whereas a crop is labeled edited only if it contains reference edit pixels. This addresses a training ambiguity in which a crop from an edited source image can itself contain no annotated edit.

The contributions of this study are as follows:
\begin{enumerate}
\item We describe a single-image detection and localization pipeline combining pretrained content features, residual filtering, cross-attention fusion, and overlapping native-resolution inference.
\item We formulate crop-aware supervision and compare a joint-view model with earlier checkpoints, examining false alarms and missed edits separately.
\item We report development results across five editor groups and edit-size ranges, supported by qualitative examples and comparisons with published methods. We distinguish the recorded 10K summary from the independently reproducible development evaluation and describe the evidence needed for a matched benchmark.
\end{enumerate}
The remainder of the paper reviews related work, presents the architecture and training procedure, and discusses detection, localization, and the limitations of the available evaluation.
